\section{Geplante Arbeitspakte und angestrebte Ergebnisse}
\label{sec:geplante-arbeitspakte-und-angestrebte-ergebnisse}

\begin{description}
    \item[Grundlagen recherchieren] Zur Vorbereitung und Aufarbeitung des Themas, was \insb in das Grundlagenkapitel einfließt, wird eine Literaturrecherche durchgeführt.
    \item[Thema darstellen] Das Einleitungskapitel, welches die Motivation, die Problemstellung und Zielsetzung der Arbeit umfasst, wird formuliert.
    \item[Ansatz zur Problemlösung ermitteln] Es wird \uA ein geeignetes Verfahren zur Regelformalisierung samt Programmiersprache ermittelt und einhergehende Implikationen benannt.
    \item[Weitere Anforderungen definieren] Es werden weitere Anforderungen, \insb an quelloffener Software, definiert.
    \item[Schnittstellen definieren] Die Schnittstelle zum Frontend wird erarbeitet, \ggfs mit dem Ersteller des Frontends diskutiert, und in einem Schnittstellenvertrag festgehalten.
    \item[Regelwerk interpretieren] Anhand des ausgewählten Verfahrens wird das Basis-Regelwerk interpretiert und konzeptionell formalisiert, sodass Anweisungen, \zB für die Implementierung eines Merkmals, entstehen.
    \item[Projektstruktur und Architektur festlegen] Vorbereitend zur Konzeptumsetzung wird die grundlegende Projektstruktur und Architektur festgelegt, umgesetzt und dokumentiert.
    \item[Design-Vorgaben festlegen] Um am Projekt mitzuwirken (dies schließt die initiale Konzeptumsetzung ein), werden \uA der Code-Style, Namensgebungen, Testabdeckung, Dokumentationspflichten und \ggfs vom Standard, \zB der Programmiersprache oder gängigen ingenieurmäßigen Software-Entwicklung, abweichende Vorgaben definiert und dokumentiert.
    Teils werden elementare, aber dennoch wichtige Vorgaben explizit festgehalten.
    \item[Konzept umsetzen] Beinhaltet die letztliche Implementierung der Kern-Funktionalitäten, wie die Frontend-Schnittstelle, Charakter-Erstellung und Regelwerks-Umsetzung.
    \item[Qualität sicherstellen] Gemäß der Design-Vorgaben muss die Qualität der Software sichergestellt werden und soll zum Projektende abschließend kontrolliert und \ggfs verbessert werden.
\end{description}
